\section{Appendix D: Proof of Theorem~\ref{thm:matrix-distribution-complete}}
\label{sec:proof-thm:matrix-distribution-complete}




As before (see \eqref{eq:assumption-lambda-r-positive}), we assume throughout the proof that $\lambda_r^{\star}>0$ for the purpose of simplifying notation. 




\subsection{Proof outline}


We now outline the proof of our distributional guarantees in Theorem~\ref{thm:matrix-distribution-complete}. 
The first step consists of justifying the heuristic first-order approximation in \eqref{eq:approx-ULambdaU-Mstar-discuss}. This is stated in the lemma below, with the proof deferred to Section~\ref{sec:sub-proof-thm:matrix-distribution}. 
%
\begin{lemma}[First-order approximation]
\label{thm:matrix-distribution}
%
Suppose that the assumptions of Theorem~\ref{thm:UsgnH-Ustar-MUstar-general} hold. Then with probability at least $1-O(n^{-5})$, one can write
%
\begin{subequations}
\label{eq:U-M-decomposition-thm}
\begin{align}
\bm{U}\mathsf{sgn}(\bm{H})-\bm{U}^{\star} & =\underset{\eqqcolon\,\bm{Z}}{\underbrace{\bm{E}\bm{U}^{\star}\big(\bm{\Lambda}^{\star}\big)^{-1}}}+\bm{\Psi},
	\label{eq:U-Ustar-decomposition-Z-Psi-thm}\\
\bm{M}-\bm{M}^{\star} & =\underset{\eqqcolon\,\bm{W}}{\underbrace{\bm{E}\bm{U}^{\star}\bm{U}^{\star\top}+\bm{U}^{\star}\bm{U}^{\star\top}\bm{E}}}+\bm{\Phi}
	\label{eq:M-Mstar-diff-decomposition-thm}
\end{align}
\end{subequations}
%
for some matrices $\bm{\Psi}$ and $\bm{\Phi}$ obeying
%
\begin{subequations}
\label{eq:Psi-Phi-bound-thm}
\begin{align}
\|\bm{\Psi}\|_{2,\infty} & \lesssim\frac{\sigma^{2}\kappa\sqrt{\mu rn\log^{2}n}}{(\lambda_{r}^{\star})^{2}}+\frac{\kappa\sigma^{2}\sqrt{\mu rn}}{(\lambda_{r}^{\star})^{2}}+\frac{\sigma}{\lambda_{r}^{\star}}\sqrt{\frac{\mu r^{2}\log n}{n}},\\
\big\|\bm{\Phi}\big\|_{\infty} & \lesssim\frac{\sigma^{2}\mu r\kappa^{2}\log n}{\lambda_{r}^{\star}}+\frac{\sigma\kappa\mu\sqrt{r^{3}\log n}}{n}.
\end{align}
\end{subequations}
%
\end{lemma}
%
\begin{remark}
	In addition to quantifying the goodness of the approximation \eqref{eq:M-Mstar-diff-decomposition-thm}, Lemma~\ref{thm:matrix-distribution} also delivers a more refined characterization for the first-order approximation $\bm{U}\mathsf{sgn}(\bm{H})-\bm{U}^{\star}\approx \bm{E}\bm{U}^{\star}\big(\bm{\Lambda}^{\star}\big)^{-1}$ in comparison to Theorem~\ref{thm:UsgnH-Ustar-MUstar-general}.  As it turns out, this result \eqref{eq:U-Ustar-decomposition-Z-Psi-thm} also assists in performing statistical inference on the low-rank factors $\bm{U}^{\star}$. The interested reader is referred  to \citet{yan2021inference} for details. 
\end{remark}


In turn, Lemma~\ref{thm:matrix-distribution} motivates one to pin down the distribution of the matrix $\bm{W}$ in \eqref{eq:M-Mstar-diff-decomposition-thm}. This can be accomplished by invoking the Berry-Esseen Theorem (e.g., \citet[Theorem 3.7]{chen2010normal}), which gives rise to the following distributional characterization. The proof of this lemma can be found in Section~\ref{sec:proof-lem:normality-spectral}. 
%
\begin{lemma}[Gaussian approximation]
	\label{lem:normality-spectral}
	Suppose that the assumptions of Theorem~\ref{thm:UsgnH-Ustar-MUstar-general} hold, and that 
	%
	%assume that
%
\begin{align}
	\label{eq:Uj-Ui-lower-bound-lem-2}
	\frac{\big\|\bm{U}_{j,\cdot}^{\star}\big\|_{2}^{2}+\big\|\bm{U}_{i,\cdot}^{\star}\big\|_{2}^{2}}{\big\|\bm{U}^{\star}\big\|_{\mathrm{F}}^{2}}
	& \gtrsim \frac{B^{2}\kappa^2 \mu^{2}r^2\log ^2 n}{\sigma_{\min}^{2}n^{2}}
	+ \frac{\sigma^{4}\mu^{2}r\kappa^{4}\log^{3}n}{\sigma_{\min}^{2}(\lambda_{r}^{\star})^{2}} .
\end{align}
%
%Condition~\ref{eq:Uj-Ui-lower-bound-lem-2} holds. 
%
	Let $\bm{W}=\bm{E}\bm{U}^{\star}\bm{U}^{\star\top} + \bm{U}^{\star}\bm{U}^{\star\top}\bm{E}$. For any $1\leq i,j\leq n$, one has
	%
	\begin{subequations}
	\begin{align}
		\sup_{z\in\mathbb{R}}\left|\mathbb{P}\left( W_{i,j}\leq z \sqrt{v_{i,j}^{\star}}\right) -\Phi(z)\right| &=o(1) , \label{eq:eqn-normality-spectral}\\
		\big\|\bm{\Phi}\big\|_{\infty} &= o\left( \sqrt{ v_{i,j}^{\star} } \right), 
		\label{eq:Phi-inf-norm-o-vij}
	\end{align}
	\end{subequations}
	%
	where $v_{i,j}^{\star}$ is defined in \eqref{eq:variance-formula-Mij-lem}, and $\Phi(\cdot)$ denotes the CDF of the standard Gaussian distribution.  
\end{lemma}
%



To finish up, invoke Lemma~\ref{thm:matrix-distribution} and \eqref{eq:Phi-inf-norm-o-vij} to yield
%
\[
	M_{i,j}-M_{i,j}^{\star}=W_{i,j}+\delta_{\phi}\sqrt{v_{i,j}^{\star}}\qquad\text{with }\delta_{\phi}=o(1).
\]
%
A little algebra further gives
%
\begin{align*}
 & \left|\mathbb{P}\left( M_{i,j}-M_{i,j}^{\star}\leq z\sqrt{v_{i,j}^{\star}}\right) -\Phi(z)\right|=\left|\mathbb{P}\left( W_{i,j}\leq\left(z-\delta_{\phi}\right)\sqrt{v_{i,j}^{\star}}\right) -\Phi(z)\right|\\
 & \qquad \leq \left|\mathbb{P}\left( W_{i,j}\leq\left(z-\delta_{\phi}\right)\sqrt{v_{i,j}^{\star}}\right)  -\Phi(z-\delta_{\phi})\right|+\left|\Phi(z-\delta_{\phi})-\Phi(z)\right|\\
 & \qquad\leq o(1)+\delta_{\phi}=o(1),
\end{align*}
%
where the last line invokes Lemma~\ref{lem:normality-spectral} and the fact that $|\Phi(u)-\Phi(v)|\leq |u-v|$ for any $u,v\in \mathbb{R}$. 
This completes the proof of  Theorem~\ref{thm:matrix-distribution-complete}, as long as Lemmas~\ref{thm:matrix-distribution} and \ref{lem:normality-spectral} can be established. 
The rest of this section is thus devoted to proving Lemmas~\ref{thm:matrix-distribution} and \ref{lem:normality-spectral}. 






\subsection{Proof of Lemma~\ref{thm:matrix-distribution}}
\label{sec:sub-proof-thm:matrix-distribution}



Before proceeding to the proof, we make note of several preliminary facts that are all direct consequences of the analysis of Theorem~\ref{thm:UsgnH-Ustar-MUstar-general} and Corollary~\ref{cor:entrywise-error-general}. The proof of these preliminary results can be found in Section~\ref{sec:proof-auxiliary-distribution}.
%
\begin{lemma}
\label{lem:bound-Delta1-Delta2}
%
With probability exceeding $1-O(n^{-5})$, one can write
%
\begin{subequations}\label{eq:decomp_soda}
\begin{align}
\bm{U}\bm{\Lambda}\mathsf{sgn}(\bm{H}) & =\bm{M}\bm{U}^{\star}+\bm{\Delta}_{1} ,
			\label{eq:U-Lambda-sgnH-decompose}\\
\mathsf{sgn}(\bm{H})\bm{\Lambda}^{\star}\mathsf{sgn}(\bm{H})^{\top} & =\bm{\Lambda}+\bm{\Delta}_{2}
			\label{eq:sgnH-Lambdastar-decompose}
\end{align}
\end{subequations}
%
for some matrices $\bm{\Delta}_{1}$ and $\bm{\Delta}_{2}$ obeying
%
\begin{subequations}
\begin{align}
  \big\| \bm{\Delta}_1 \big\|_{2,\infty} & \lesssim 
\frac{\sigma^{2}\kappa\sqrt{\mu rn\log^{2}n}}{\lambda_{r}^{\star}}, \label{eq:Delta-1-U-R-Lambda}\\
  \big\|\bm{\Delta}_2\big\|
%\\
% & \qquad\leq\big\|\mathsf{sgn}(\bm{H})\bm{\Lambda}^{\star}\mathsf{sgn}(\bm{H})^{\top}-\bm{H}\bm{\Lambda}^{\star}\bm{H}^{\top}\big\|+\big\|\bm{H}\bm{\Lambda}^{\star}\bm{H}^{\top}-\bm{\Lambda}\big\|\\
 %& \qquad
  &\lesssim\frac{\kappa\sigma^{2}n}{\lambda_{r}^{\star}}+\sigma\sqrt{r\log n}.
   \label{eq:sgn-H-Lambdastar-sgn-H-Lambda}
\end{align}
\end{subequations}
%
\end{lemma}
%



We are now ready to embark on the proof of Lemma~\ref{thm:matrix-distribution}. 
In order to analyze the behavior of $\bm{U}\mathsf{sgn}(\bm{H})$, 
we first point out the following decomposition: 
%
\begin{align*}
\bm{U}\mathsf{sgn}(\bm{H})\bm{\Lambda}^{\star} & =\bm{U}\bm{\Lambda}\mathsf{sgn}(\bm{H})+\bm{U}\big(\mathsf{sgn}(\bm{H})\bm{\Lambda}^{\star}-\bm{\Lambda}\mathsf{sgn}(\bm{H})\big)\\
 & =\bm{M}\bm{U}^{\star}+\bm{\Delta}_{1}+\bm{U}\big(\mathsf{sgn}(\bm{H})\bm{\Lambda}^{\star}\mathsf{sgn}(\bm{H})^{\top}-\bm{\Lambda}\big)\mathsf{sgn}(\bm{H})\\
 & =\bm{M}^{\star}\bm{U}^{\star}+\bm{E}\bm{U}^{\star}+\bm{\Delta}_{1}+\bm{U}\bm{\Delta}_{2} \, \mathsf{sgn}(\bm{H})\\
 & =\bm{U}^{\star}\bm{\Lambda}^{\star}+\bm{E}\bm{U}^{\star}+\bm{\Delta}_{1}+\bm{U}\bm{\Delta}_{2}\, \mathsf{sgn}(\bm{H}),
\end{align*}
%
where the second and the third identities result from Lemma~\ref{lem:bound-Delta1-Delta2} (cf.~\eqref{eq:decomp_soda}). 
The key point of this decomposition is to establish a connection between $\bm{U}\mathsf{sgn}(\bm{H})\bm{\Lambda}^{\star}$ and $\bm{U}^{\star}\bm{\Lambda}^{\star} + \bm{E}\bm{U}^{\star}$, with the assistance of the matrices $\bm{\Delta}_1$ and $\bm{\Delta}_2$ studied in Lemma~\ref{lem:bound-Delta1-Delta2}. 
A little algebra then yields
%
\begin{align}
\bm{U}\mathsf{sgn}(\bm{H})-\bm{U}^{\star}=\underset{\eqqcolon\,\bm{Z}}{\underbrace{\bm{E}\bm{U}^{\star}\big(\bm{\Lambda}^{\star}\big)^{-1}}} & +\underset{\eqqcolon\,\bm{\Psi}}{\underbrace{\bm{\Delta}_{1}\big(\bm{\Lambda}^{\star}\big)^{-1}+\bm{U}\bm{\Delta}_{2}\,\mathsf{sgn}(\bm{H})\big(\bm{\Lambda}^{\star}\big)^{-1}}}.
\label{eq:U-Ustar-decomposition-Z-Psi}
\end{align}
%
In view of Lemma~\ref{lem:bound-Delta1-Delta2}, the residual matrix
$\bm{\Psi}$ obeys
%
\begin{align}
\|\bm{\Psi}\|_{2,\infty} & \leq\big\|\bm{\Delta}_{1}\big\|_{2,\infty}\big\|\big(\bm{\Lambda}^{\star}\big)^{-1}\big\|+\big\|\bm{U}\big\|_{2,\infty}\big\|\bm{\Delta}_{2}\big\|\,\big\|\mathsf{sgn}(\bm{H})\big\|\,\big\|\big(\bm{\Lambda}^{\star}\big)^{-1}\big\|\nonumber \\
 & \lesssim\frac{1}{\lambda_{r}^{\star}}\big\|\bm{\Delta}_{1}\big\|_{2,\infty}+\frac{1}{\lambda_{r}^{\star}}\sqrt{\frac{\mu r}{n}}\big\|\bm{\Delta}_{2}\big\|\nonumber \\
 & \lesssim\frac{\sigma^{2}\kappa\sqrt{\mu rn\log^{2}n}}{(\lambda_{r}^{\star})^{2}}+\frac{\kappa\sigma^{2}\sqrt{\mu rn}}{(\lambda_{r}^{\star})^{2}}+\frac{\sigma}{\lambda_{r}^{\star}}\sqrt{\frac{\mu r^{2}\log n}{n}}
\label{eq:Psi-two-inf-norm-bound}
\end{align}
%
as claimed. 



The next step lies in analyzing the matrix estimator $\bm{M}=\bm{U}\bm{\Lambda}\bm{U}^{\top}$.
Towards this, we make the observation that
%
\begin{align}
\bm{M}-\bm{M}^{\star} & =\bm{U}\bm{\Lambda}\bm{U}^{\top}-\bm{M}^{\star}\nonumber \\
 & =\bm{U}\big(\mathsf{sgn}(\bm{H})\bm{\Lambda}^{\star}\mathsf{sgn}(\bm{H})^{\top}\big)\bm{U}^{\top}-\bm{U}\bm{\Delta}_{2}\bm{U}^{\top}-\bm{M}^{\star}\nonumber \\
 & =(\bm{U}^{\star}+\bm{Z}+\bm{\Psi})\bm{\Lambda}^{\star}(\bm{U}^{\star}+\bm{Z}+\bm{\Psi})^{\top}-\bm{U}\bm{\Delta}_{2}\bm{U}^{\top}-\bm{U}^{\star}\bm{\Lambda}^{\star}\bm{U}^{\star\top}\nonumber \\
 & =\bm{Z}\bm{\Lambda}^{\star}\bm{U}^{\star\top}+\bm{U}^{\star}\bm{\Lambda}^{\star}\bm{Z}^{\top}+\bm{\Phi}\nonumber \\
 & =\bm{E}\bm{U}^{\star}\bm{U}^{\star\top}+\bm{U}^{\star}\bm{U}^{\star\top}\bm{E}+\bm{\Phi},\label{eq:M-Mstar-diff-decomposition}
\end{align}
%
where the second line relies on Lemma~\ref{lem:bound-Delta1-Delta2} (cf.~\eqref{eq:sgnH-Lambdastar-decompose}),
the third identity makes use of (\ref{eq:U-Ustar-decomposition-Z-Psi}),
and the residual matrix $\bm{\Phi}$ is defined as
%
\begin{equation}
\bm{\Phi}\coloneqq\bm{\Psi}\bm{\Lambda}^{\star}\big(\bm{U}\mathsf{sgn}(\bm{H})\big)^{\top}+\bm{U}^{\star}\bm{\Lambda}^{\star}\bm{\Psi}^{\top}+\bm{Z}\bm{\Lambda}^{\star}\big(\bm{U}\mathsf{sgn}(\bm{H})-\bm{U}^{\star}\big)^{\top}-\bm{U}\bm{\Delta}_{2}\bm{U}^{\top}.\label{eq:defn-Phi-residual}
\end{equation}
%
In addition, it is seen from (\ref{eq:EU-2-inf-bound-general}) that
%
\begin{equation}
\|\bm{Z}\|_{2,\infty}\leq\big\|\bm{E}\bm{U}^{\star}\big(\bm{\Lambda}^{\star}\big)^{-1}\big\|_{2,\infty}\leq\frac{1}{\lambda_{r}^{\star}}\big\|\bm{E}\bm{U}^{\star}\big\|_{2,\infty}\lesssim\frac{\sigma\sqrt{r\log n}}{\lambda_{r}^{\star}}.\label{eq:Z-two-inf-norm-bound}
\end{equation}
%
Consequently, one can deduce that
%
\begin{align*}
\big\|\bm{\Phi}\big\|_{\infty} & \leq\big\|\bm{\Lambda}^{\star}\big\|\left\{ \|\bm{Z}\|_{2,\infty}^{2}+\|\bm{Z}\|_{2,\infty}\big\|\bm{U}\mathsf{sgn}(\bm{H})-\bm{U}^{\star}\big\|_{2,\infty}\right\} \\
 & \quad+\big\|\bm{\Lambda}^{\star}\big\|\|\bm{\Psi}\|_{2,\infty}\left(\big\|\bm{U}\big\|_{2,\infty}+\|\bm{U}^{\star}\|_{2,\infty}\right)+\|\bm{U}\|_{2,\infty}^{2}\|\bm{\Delta}_{2}\|\\
 & \lesssim\big|\lambda_{1}^{\star}\big|\left\{ \frac{\sigma^{2}r\log n}{(\lambda_{r}^{\star})^{2}}+\frac{\sigma\sqrt{r\log n}}{\lambda_{r}^{\star}}\cdot\frac{\sigma\kappa\sqrt{\mu r\log n}}{\lambda_{r}^{\star}}\right\} \\
 & \quad+\big|\lambda_{1}^{\star}\big|\sqrt{\frac{\mu r}{n}}\left\{ \frac{\sigma^{2}\kappa\sqrt{\mu rn\log^{2}n}}{(\lambda_{r}^{\star})^{2}}+\frac{\kappa\sigma^{2}\sqrt{\mu rn}}{(\lambda_{r}^{\star})^{2}}+\frac{\sigma}{\lambda_{r}^{\star}}\sqrt{\frac{\mu r^{2}\log n}{n}}\right\} \\
 & \quad+\frac{\mu r}{n}\left\{ \frac{\kappa\sigma^{2}n}{\lambda_{r}^{\star}}+\sigma\sqrt{r\log n}\right\} \\
 & \lesssim\frac{\sigma^{2}\mu r\kappa^{2}\log n}{\lambda_{r}^{\star}}+\frac{\sigma\kappa\mu\sqrt{r^{3}\log n}}{n},
\end{align*}
%
where the second inequality follows from (\ref{eq:Psi-two-inf-norm-bound}),
(\ref{eq:Z-two-inf-norm-bound}), \eqref{eq:incoherence-U-estimate}, Theorem~\ref{thm:UsgnH-Ustar-MUstar-general}, and Lemma~\ref{lem:bound-Delta1-Delta2}. 

%






\subsection{Proof of Lemma~\ref{lem:normality-spectral}}
\label{sec:proof-lem:normality-spectral}

In what follows, we shall only focus on the case with $i\neq j$. The case with $i=j$ can be analyzed in an analogous manner; we omit it for the sake of brevity.
Before proceeding to the proof, we make note of a couple of basic facts about $\bm{P}^{\star}$ that will prove useful.  
The first property asserts that, for any $1\leq j\leq n$, 
%
\begin{align}
\sum_{l=1}^{n}P_{j,l}^{\star2} & =\sum_{l=1}^{n}P_{l,j}^{\star2}=\big\|\big(\bm{U}^{\star}\bm{U}^{\star\top}\big)_{\cdot,j}\big\|_{2}^{2}=\big\|\bm{U}^{\star}\big(\bm{U}_{j,\cdot}^{\star}\big)^{\top}\big\|_{2}^{2}=\bm{U}_{j,\cdot}^{\star}\bm{U}^{\star\top}\bm{U}^{\star}\big(\bm{U}_{j,\cdot}^{\star}\big)^{\top}\nonumber \\
 & =\big\|\bm{U}_{j,\cdot}^{\star}\big\|_{2}^{2}=P_{j,j}^{\star}.
	\label{eq:P-jl-square-sum}
\end{align}
%
The second property is concerned with the term $\|\bm{P}^{\star}\|_{\infty}$:  
%
\begin{equation}
	\|\bm{P}^{\star}\|_{\infty}=\|\bm{U}^{\star}\bm{U}^{\star\top}\|_{\infty}\leq\|\bm{U}^{\star}\|_{2,\infty}^{2}\leq\frac{\mu r}{n},
	\label{eq:P-inf-norm-UB}
\end{equation}
%
where the last inequality follows from the incoherence assumption. 




In view of the definition \eqref{eq:defn-P-UU-T} of $\bm{P}^{\star}$, we can express $\bm{W}=\bm{E}\bm{P}^{\star}+\bm{P}^{\star}\bm{E}$, which reveals that 
%
\begin{equation}
	W_{i,j}
	%=\sum_{l=1}^{n}E_{i,l}P_{l,j}+\sum_{l=1}^{n}P_{i,l}E_{l,j}
	=\sum_{l:\,l\neq j}E_{i,l}P^{\star}_{l,j}+\sum_{l:\,l\neq i}P^{\star}_{i,l}E_{l,j}+E_{i,j}(P^{\star}_{i,i}+P^{\star}_{j,j}).
\end{equation}
%
In other words,  $W_{i,j}$ can be viewed as a weighted sum of independent random variables $\{E_{i,l}\mid l\neq j\} \cup \{E_{l,j}\mid l\neq i\} \cup \{E_{i,j}\}$. 
To pin down the distribution of $W_{i,j}$, we resort to a non-asymptotic
version of the celebrated Berry-Esseen Theorem; 
see \citet[Theorem 3.7]{chen2010normal} for a proof using Stein's method.
%
\begin{theorem}[The Berry-Esseen bound]\label{thm:berry-esseen}
Let $\xi_{1},\ldots,\xi_{n}$ be independent zero-mean random variables satisfying $\sum_{i=1}^{n}\mathsf{Var}(\xi_{i})= v$.
Then the quantity $S= \frac{1}{\sqrt{v}} \sum_{i=1}^{n}\xi_{i}$ satisfies 
%
\[
	\sup_{z\in\mathbb{R}}\big|\,\mathbb{P}\left(S\leq z\right)-\Phi\left(z\right)\big|\leq10\gamma,\qquad\text{where}\ \gamma
	=\sum_{i=1}^{n} \frac{\mathbb{E}\big[|\xi_{i}|^{3}\big] }{ v^{3/2} }.
\]
%
%where $\Phi(\cdot)$ represents the CDF of a standard Gaussian distribution. 
\end{theorem}





According to the Berry-Esseen bound (cf.~Theorem~\ref{thm:berry-esseen}), 
proving the approximate Gaussianity of $W_{i,j}$ boils down to characterizing the second and the third moments of these random variables under consideration. 




Let us start with the variance statistics. 
Given that $\{E_{i,j}\mid i\geq j\}$ are independently generated, 
%the variance  $v_{i,j}^{\star}  \coloneqq\mathsf{Var}(W_{i,j}) $ 
we can straightforwardly see that
%
\begin{align}
	v_{i,j}^{\star}   %\coloneqq\mathsf{Var}(W_{i,j})
	&=\sum_{l:\,l\neq j}\sigma_{i,l}^{2}P_{l,j}^{\star 2}+\sum_{l:\,l\neq i}P_{i,l}^{\star 2}\sigma_{l,j}^{2}+\sigma_{i,j}^{2}(P_{i,i}^{\star}+P_{j,j}^{\star})^{2} \notag\\
	&= \mathsf{Var}(W_{i,j}).
	\label{eq:variance-formula-Mij}
\end{align}
%
We now develop a lower bound on this variance term.  Given
that $\bm{P}^{\star} \succeq \bm{0}$, one has $P_{i,i}^{\star},P_{j,j}^{\star}\geq0$,
%and $(P^{\star}_{i,i}+P^{\star}_{j,j})^{2}\geq P_{i,i}^{\star 2}+P_{j,j}^{\star 2}$,
which combined with \eqref{eq:P-jl-square-sum} reveals that
%
\begin{align}
v_{i,j}^{\star} & 
%\geq\sigma_{\min}^{2}\Bigg\{\sum_{l:\,l\neq i}P_{l,j}^{\star 2}+\sum_{l:\,l\neq j}P_{i,l}^{\star 2}+P_{i,i}^{\star 2}+P_{j,j}^{\star 2}\Bigg\}
\geq \sigma_{\min}^{2}\Bigg\{\sum_{l=1}^{n}P_{l,j}^{\star 2}+\sum_{l=1}^{n}P_{i,l}^{\star 2}\Bigg\} 
%\notag\\
% & =\sigma_{\min}^{2}\Big\{\big\|\big(\bm{U}^{\star}\bm{U}^{\star\top}\big)_{\cdot,j}\big\|_{2}^{2}+\big\|\big(\bm{U}^{\star}\bm{U}^{\star\top}\big)_{i,\cdot}\big\|_{2}^{2}\Big\} \notag\\
% & =\sigma_{\min}^{2}\Big\{\big\|\bm{U}^{\star}\big(\bm{U}_{j,\cdot}^{\star}\big)^{\top}\big\|_{2}^{2}+\big\|\bm{U}_{i,\cdot}^{\star}\bm{U}^{\star\top}\big\|_{2}^{2}\big\} \notag\\
% & 
=\sigma_{\min}^{2}\Big\{\big\|\bm{U}_{j,\cdot}^{\star}\big\|_{2}^{2}+\big\|\bm{U}_{i,\cdot}^{\star}\big\|_{2}^{2}\Big\}.
	\label{eq:vij-lower-bound-12}
\end{align}
%

Next, we move on to bound the third moments. 
Utilizing the independence of $\{E_{i,j}\mid i\geq j\}$ once again gives
%
\begin{align}
	\gamma & \coloneqq\frac{\sum_{l:\,l\neq j}\mathbb{E}\big[\big|E_{i,l}\big|^{3}\big]\big|P_{l,j}^{\star} \big|^{3}+\sum_{l:\,l\neq i}\big|P_{i,l}\big|^{3}\mathbb{E}\big[\big|E_{l,j}\big|^{3}\big]+\mathbb{E}\big[\big|E_{i,j}\big|^{3}\big]\big|P_{i,i}^{\star} +P_{j,j}^{\star} \big|^{3}}{\big(v_{i,j}^{\star}\big)^{3/2}} \notag\\
 & \leq\frac{2B\|\bm{P}^{\star} \|_{\infty}}{\big(v_{i,j}^{\star}\big)^{3/2}}\left\{ \sum_{l:\,l\neq j}\mathbb{E}\big[E_{i,l}^{2}\big] P_{l,j}^{\star 2}+\sum_{l:\,l\neq i} P_{i,l}^{\star 2}\mathbb{E}\big[E_{l,j}^{2}\big]+\mathbb{E}\big[E_{i,j}^{2}\big]\big|P_{i,i}^{\star} +P_{j,j}^{\star} \big|^{2}\right\} \notag\\
 & =\frac{2B\|\bm{P}^{\star} \|_{\infty}}{\big(v_{i,j}^{\star}\big)^{3/2}}\cdot v_{i,j}^{\star}=\frac{2B\|\bm{P}^{\star} \|_{\infty}}{\big(v_{i,j}^{\star}\big)^{1/2}}, 
	\label{eq:gamma-UB-123}
\end{align}
% 
where we have used the assumption that $|E_{i,j}|\leq B$, and the last line arises from the expression \eqref{eq:variance-formula-Mij}. 
Substituting \eqref{eq:P-inf-norm-UB} and \eqref{eq:vij-lower-bound-12} into \eqref{eq:gamma-UB-123} and invoking the elementary identity $\|\bm{U}^{\star} \|_{\mathrm{F}}^2=r$, 
we arrive at
%
\begin{align*}
\gamma & \leq\frac{2B\mu r}{n\sigma_{\min}\sqrt{\big\|\bm{U}_{j,\cdot}^{\star}\big\|_{2}^{2}+\big\|\bm{U}_{i,\cdot}^{\star}\big\|_{2}^{2}}}=\frac{2B\mu\sqrt{r}\,\big\|\bm{U}^{\star}\big\|_{\mathrm{F}}}{n\sigma_{\min}\sqrt{\big\|\bm{U}_{j,\cdot}^{\star}\big\|_{2}^{2}+\big\|\bm{U}_{i,\cdot}^{\star}\big\|_{2}^{2}}}
	= o(1),
\end{align*}
%
provided that Condition~\eqref{eq:Uj-Ui-lower-bound-lem-2} holds. 


With the above calculations in place, invoking Theorem~\ref{thm:berry-esseen} immediately leads to 
%
\[
\sup_{z\in\mathbb{R}}\left|\mathbb{P}\left( W_{i,j}\leq\sqrt{v_{i,j}^{\star}}z\right) -\Phi(z)\right|\leq10\gamma=o(1). 
\]
%
as claimed in \eqref{eq:eqn-normality-spectral}. 



Finally, we turn to proving the bound \eqref{eq:Phi-inf-norm-o-vij}. 
By virtue of Lemma~\ref{thm:matrix-distribution} and \eqref{eq:vij-lower-bound-12}, we know that 
%
\begin{align*}
\big\|\bm{\Phi}\big\|_{\infty} & \lesssim\frac{\sigma^{2}\mu r\kappa^{2}\log n}{\lambda_{r}^{\star}}+\frac{\sigma\kappa\mu\sqrt{r^{3}\log n}}{n}\\
 & =o\left(\sigma_{\min}\sqrt{\big\|\bm{U}_{j,\cdot}^{\star}\big\|_{2}^{2}+\big\|\bm{U}_{i,\cdot}^{\star}\big\|_{2}^{2}}\right)\leq o\big(\sqrt{v_{i,j}^{\star}}\big) , 
\end{align*}
%
with the proviso that
%
\begin{align}
	& \sqrt{\big\|\bm{U}_{j,\cdot}^{\star}\big\|_{2}^{2}+\big\|\bm{U}_{i,\cdot}^{\star}\big\|_{2}^{2}}  \gtrsim\frac{\sigma^{2}\mu r\kappa^{2}\log^{3/2}n}{\sigma_{\min}\lambda_{r}^{\star}}+\frac{\sigma\kappa\mu\sqrt{r^{3}}\log n}{\sigma_{\min}n} \notag\\
 & \qquad =\left(\frac{\sigma^{2}\mu\sqrt{r}\kappa^{2}\log^{3/2}n}{\sigma_{\min}\lambda_{r}^{\star}}+\frac{\sigma\kappa\mu r\log n}{\sigma_{\min}n}\right)\big\|\bm{U}^{\star}\big\|_{\mathrm{F}}. 
	\label{eq:condition-7878}
\end{align}
%
Recognizing the trivial bound $\sigma^2 = \max_{i,j}\mathbb{E}[E_{i,j}^2]\leq B^2$, we know that Condition~\eqref{eq:condition-7878} holds as long as
%
\begin{align}
	\frac{\big\|\bm{U}_{j,\cdot}^{\star}\big\|_{2}^{2}+\big\|\bm{U}_{i,\cdot}^{\star}\big\|_{2}^{2}}{\big\|\bm{U}^{\star}\big\|_{\mathrm{F}}^{2}}
	& \gtrsim \frac{\sigma^{4}\mu^{2}r\kappa^{4}\log^{3}n}{\sigma_{\min}^{2}(\lambda_{r}^{\star})^{2}} + \frac{B^{2}\kappa^2 \mu^{2}r^2\log ^2 n}{\sigma_{\min}^{2}n^{2}}, 
	%+\frac{\sigma^{2}\kappa^{2}\mu^{2}r^{2}\log^{2}n}{\sigma_{\min}^{2}n^{2}}.
	%\label{eq:Uj-Ui-lower-bound-lem-2}
\end{align}
%
which is precisely Condition~\eqref{eq:Uj-Ui-lower-bound-lem-2}. This concludes the proof of Lemma~\ref{lem:normality-spectral}. 

%as claimed. 


%
%\[
%\frac{B\|\bm{P}\|_{\infty}}{\big(v_{i,j}^{\star}\big)^{1/2}}\lesssim\frac{\sigma_{\max}\sqrt{n/(\mu^{2}\log n)}\cdot\frac{\mu r}{n}}{\sigma_{\min}\Big\{\Big\|\bm{U}_{j,\cdot}^{\star}\Big\|_{2}+\Big\|\bm{U}_{i,\cdot}^{\star}\Big\|_{2}\Big\}}\asymp\frac{1}{\sqrt{\log n}}\frac{\|\bm{U}^{\star}\|_{\mathrm{F}}}{\sqrt{n}\Big\{\big\|\bm{U}_{j,\cdot}^{\star}\big\|_{2}+\big\|\bm{U}_{i,\cdot}^{\star}\big\|_{2}\Big\}}
%\]
%


%And we know that 
%\[
%B\lesssim\sigma_{\max}\sqrt{n/(\mu^{2}\log n)}.
%\]






\subsection{Proof of Lemma~\ref{lem:bound-Delta1-Delta2}}
\label{sec:proof-auxiliary-distribution}



%\paragraph{Proof of Lemma~\ref{lem:bound-Delta1-Delta2}.}
To begin with, let us begin by proving \eqref{eq:Delta-1-U-R-Lambda}. 
From the definition of the quantity $\mathcal{E}_{1}$ (see Lemma~\ref{lem:UH-MUstar-diff-decompose}), we have
%
\begin{align*}
 & \big\|\bm{M}(\bm{U}\bm{H}-\bm{U}^{\star})\big\|_{2,\infty}=\lambda_{r}^{\star}\mathcal{E}_{1}/2\\
 & \qquad\lesssim\left(\alpha_{0}+\alpha_{1}+\alpha_{2}\right)+\left(\sigma\sqrt{n}+B\log n\right)\big\|\bm{U}\bm{H}-\bm{U}^{\star}\big\|_{2,\infty}\\
 & \qquad\lesssim\left(\alpha_{0}+\alpha_{1}+\alpha_{2}\right)+\left(\sigma\sqrt{n}+B\log n\right)\frac{\sigma\kappa\sqrt{\mu r\log n}}{\lambda_{r}^{\star}}\\
 & \qquad\asymp\frac{\sigma^{2}\kappa\sqrt{\mu rn\log^{2}n}}{\lambda_{r}^{\star}},
\end{align*}
%
where the first inequality comes from \eqref{eq:E1-UB-general-1} and \eqref{eq:defn-E1-rho1},
the second inequality is a consequence of Theorem~\ref{thm:UsgnH-Ustar-MUstar-general}, and 
the last line relies on our previous bounds on $\alpha_0,\alpha_1,\alpha_2$ (see \eqref{eq:defn-alpha-0-general}, \eqref{eq:alpha1-bound-1234} and \eqref{eq:Mstar-UH-Ustar-2inf}) and holds as long as $B\lesssim\sigma\sqrt{n/(\mu\log n)}$. 
%
Additionally, from the elementary identity $\bm{M}\bm{U}\bm{H}=\bm{U}\bm{\Lambda}\bm{H}$, we obtain
%
\begin{align*}
\big\|\bm{U}\bm{\Lambda}\mathsf{sgn}(\bm{H})-\bm{M}\bm{U}\bm{H}\big\|_{2,\infty} & =\big\|\bm{U}\bm{\Lambda}\mathsf{sgn}(\bm{H})-\bm{U}\bm{\Lambda}\bm{H}\big\|_{2,\infty}\\
 & \leq\big\|\bm{U}\big\|_{2,\infty}\big\|\bm{\Lambda}\big\|\,\big\|\mathsf{sgn}(\bm{H})-\bm{H}\big\|\\
 & \lesssim\sqrt{\frac{\mu r}{n}} |\lambda_{1}^{\star} | \cdot\frac{\sigma^{2}n}{(\lambda_{r}^{\star})^{2}}=\frac{\kappa\sigma^{2}\sqrt{\mu rn}}{\lambda_{r}^{\star}},
\end{align*}
%
where the last line results from Lemma~\ref{lem:H-property-summary-general}, the fact \eqref{eq:incoherence-U-estimate}, and the following inequality
%
\begin{equation}
	\|\bm{\Lambda}\| \leq \|\bm{\Lambda}^{\star}\| + \|\bm{E}\|\leq |\lambda_1^{\star}| + O(\sigma\sqrt{n}) \leq 2|\lambda_1^{\star}|. 
	\label{eq:Lambda-norm-upper-bound}
\end{equation}
%
Taking together the above bounds and applying the triangle inequality 
immediately establish~\eqref{eq:Delta-1-U-R-Lambda}. 







Next, we turn to the proof of the bound \eqref{eq:sgn-H-Lambdastar-sgn-H-Lambda}. 
Note that it has been shown in (\ref{eq:H-Lambda-star-H-Lambda-dist}) that
%
\[
\big\|\bm{H}\bm{\Lambda}^{\star}\bm{H}^{\top}-\bm{\Lambda}\big\|\lesssim\frac{\kappa\sigma^{2}n}{\lambda_{r}^{\star}}+\sigma\sqrt{r\log n}.
\]
%
In addition, the triangle inequality leads to
%
\begin{align*}
 & \big\|\mathsf{sgn}(\bm{H})\bm{\Lambda}^{\star}\mathsf{sgn}(\bm{H})^{\top}-\bm{H}\bm{\Lambda}^{\star}\bm{H}^{\top}\big\|\\
 & \qquad\leq\big\|\mathsf{sgn}(\bm{H})\bm{\Lambda}^{\star}\big(\mathsf{sgn}(\bm{H})-\bm{H}\big)^{\top}\big\|+\big\|\big(\mathsf{sgn}(\bm{H})-\bm{H}\big)\bm{\Lambda}^{\star}\bm{H}^{\top}\big\|\\
 & \qquad\leq\left(\big\|\mathsf{sgn}(\bm{H})\big\|+\big\|\bm{H}\big\|\right)\,\big\|\bm{\Lambda}^{\star}\big\|\,\big\|\mathsf{sgn}(\bm{H})-\bm{H}\big\|\\
 & \qquad\leq  2 |\lambda_{1}^{\star}| \, \big\|\mathsf{sgn}(\bm{H})-\bm{H}\big\|\\
 & \qquad\lesssim |\lambda_{1}^{\star}| \frac{\sigma^{2}n}{\big(\lambda_{r}^{\star}\big)^{2}}=\frac{\kappa\sigma^{2}n}{\lambda_{r}^{\star}},
\end{align*}
%
where the third line follows since $\big\|\mathsf{sgn}(\bm{H})\big\|=1$
and $\|\bm{H}\|\leq\|\bm{U}\|\|\bm{U}^{\star}\|=1$, and the last
inequality comes from (\ref{eq:H-sgnH-diff-UB-123}). Combining
the above two results and invoking the triangle inequality lead to
the advertised bound \eqref{eq:sgn-H-Lambdastar-sgn-H-Lambda}. 





\section{Appendix E: Proof of Theorem~\ref{thm:CI-general}} 
\label{sec:proof-thm:CI-general}


With the distributional guarantees in Theorem~\ref{thm:matrix-distribution-complete} in place, 
the only remaining task boils down to verifying the statistical accuracy of the variance estimator $\widehat{v}_{i,j}$. 
This can be achieved via the following lemma, whose proof is provided in Section~\ref{sec:proof-lem:variance-estimation-guarantee}.  
%
\begin{lemma}
	\label{lem:variance-estimation-guarantee}
	%
	Suppose that the assumptions of Theorem~\ref{thm:UsgnH-Ustar-MUstar-general} hold. In addition, assume that $\kappa^{4}\mu^{2}r^{2}\log n\leq n$, $\sigma\sqrt{n}\lesssim|\lambda_{r}^{\star}|/\kappa$ and
	%
	\begin{align}
		\big\|\bm{U}_{j,\cdot}^{\star}\big\|_{2}^{2}+\big\|\bm{U}_{i,\cdot}^{\star}\big\|_{2}^{2}\gtrsim\frac{B\sigma\kappa^{2}\mu^{2}r^{2}\log n}{\sigma_{\min}^2 n^{3/2}}+\frac{\sigma^{3}\kappa\mu^{2}r^{2}\log n}{\sigma_{\min}^2|\lambda_{r}^{\star}|\sqrt{n}}. 
		\label{eq:Uj-Ui-lower-bound-lem}
	\end{align}
	%
	With probability exceeding $1-O(n^{-5})$, one has
	%
	\begin{align}
		\big|\widehat{v}_{i,j}-v_{i,j}^{\star}\big| =o\big( v_{i,j}^{\star} \big). 
	\end{align}
	%
\end{lemma}
%

Lemma~\ref{lem:variance-estimation-guarantee} essentially enables us to express
%
\[
	\widehat{v}_{i,j} = (1 + \zeta_v)^2 v_{i,j}^{\star} \qquad \text{with } \zeta_v = o(1). 
\]
%
As a consequence, we can further demonstrate that
%
\begin{align*}
 & \left|\mathbb{P}\left( M_{i,j}^{\star}\in\mathsf{CI}_{i,j}^{1-\alpha}\right) -(1-\alpha)\right|=\left|\mathbb{P}\left( \big| \widehat{M}_{i,j}-M_{i,j}^{\star}\big|\leq z_{\alpha/2}\sqrt{\widehat{v}_{i,j}}\right) -(1-\alpha)\right|\\
 & =\left|\mathbb{P}\left( \big| \widehat{M}_{i,j}-M_{i,j}^{\star}\big|\leq(1+\zeta_{v})z_{\alpha/2}\sqrt{v_{i,j}^{\star}}\right) -(1-\alpha)\right|\\
 & \leq\left|\Phi\big((1+\zeta_{v})z_{\alpha/2}\big)-\Phi\big(-(1+\zeta_{v})z_{\alpha/2}\big)-(1-\alpha)\right|+o(1)\\
 & \leq\left|\Phi\big(z_{\alpha/2}\big)-\Phi\big(-z_{\alpha/2}\big)-(1-\alpha)\right|+2\left|\Phi\big((1+\zeta_{v})z_{\alpha/2}\big)-\Phi\big(z_{\alpha/2}\big)\right|+o(1)\\
 & \leq2\zeta_{v}z_{\alpha/2}=o(1),
\end{align*}
%
where the first inequality follows from Theorem~\ref{thm:matrix-distribution-complete} and $z_{\alpha/2}:= \Phi^{-1}(1-\alpha/2)$, the second inequality applies the triangle inequality, 
and the validity of the last line can be seen from the basic fact $|\Phi(u)-\Phi(v)| \leq |u-v|$.   


When $\sigma_{\min}\asymp \sigma$, Condition~\eqref{eq:Uj-Ui-lower-bound-lem} simplifies to
%
\begin{align}
	\frac{\big\|\bm{U}_{j,\cdot}^{\star}\big\|_{2}^{2}+\big\|\bm{U}_{i,\cdot}^{\star}\big\|_{2}^{2}}{\big\|\bm{U}^{\star}\big\|_{\mathrm{F}}^{2}}
	\gtrsim\frac{B\kappa^{2}\mu^{2}r\log n}{\sigma n^{3/2}}+\frac{\sigma\kappa\mu^{2}r\log n}{|\lambda_{r}^{\star}|\sqrt{n}}.		 
		\label{eq:Uj-Ui-lower-bound-lem-31}
\end{align}
%
We still need to ensure that Condition~\eqref{eq:Ui-Uj-condition-thm-distribution-1} is satisfied.  
It is seen that 
%
\begin{align*}
	\frac{B^{2}\kappa^{2}\mu^{2}r^{2}\log^{2}n}{\sigma^{2}n^{2}}+\frac{\sigma^{2}\mu^{2}r\kappa^{4}\log^{3}n}{(\lambda_{r}^{\star})^{2}}&\lesssim\frac{B\kappa^{2}\mu^{2}r^{2}\log^{2}n}{\sigma n^{3/2}}+\frac{\sigma\mu^{2}r\kappa^{3}\log^{3}n}{|\lambda_{r}^{\star}|\sqrt{n}},
\end{align*}
%
%where the first inequality has made use of the fact that $B\geq \sigma$, and the second inequality 
which holds if $\sigma\kappa \sqrt{n} \lesssim |\lambda_r^{\star}|$ and $B\lesssim \sigma\sqrt{n}$. 
As a result, if Condition~\eqref{eq:Uj-Ui-lower-bound-thm} holds, then both \eqref{eq:Uj-Ui-lower-bound-lem-31} and \eqref{eq:Ui-Uj-condition-thm-distribution-1}
are satisfied. 
This finishes the proof, as long as  Lemma~\ref{lem:variance-estimation-guarantee} can be established. 



%%
%\begin{align*}
% & \frac{B^{2}\mu^{2}r\log n}{n^{2}\sigma^{2}}+\frac{\sigma^{2}\mu^{2}r\kappa^{4}\log^{3}n}{(\lambda_{r}^{\star})^{2}}+\frac{\kappa^{2}\mu^{2}r^{2}\log^{2}n}{n^{2}}\\
% & \lesssim\frac{B^{2}\kappa^{2}\mu^{2}r^{2}\log^{2}n}{n^{2}\sigma^{2}}+\frac{\sigma^{2}\mu^{2}r\kappa^{4}\log^{3}n}{(\lambda_{r}^{\star})^{2}}\lesssim\frac{B\kappa^{2}\mu^{2}r^{2}\log^{2}n}{\sigma n^{3/2}}+\frac{\sigma\mu^{2}r\kappa^{3}\log^{3}n}{|\lambda_{r}^{\star}|\sqrt{n}},
%\end{align*}
%%


\subsection{Proof for Lemma~\ref{lem:variance-estimation-guarantee}}
\label{sec:proof-lem:variance-estimation-guarantee}

As before, we shall only present the proof for the case with $i\neq j$ for the sake of conciseness. 
In order to justify the goodness of the estimator $\widehat{v}_{i,j}$, 
we find it convenient to first look at the surrogate estimator introduced in \eqref{eq:v-ij-plugin-surrogate-123}, i.e., 
%
\begin{equation}
	\widetilde{v}_{i,j}=\sum_{l=1}^{n}E_{i,l}^{2}P_{l,j}^{\star 2}+\sum_{l=1}^{n}P_{i,l}^{\star 2}E_{l,j}^{2}+2E_{i,j}^{2}P_{i,i}^{\star}P_{j,j}^{\star}, 
	\label{eq:v-surrogate}
\end{equation}
%
% which can also be viewed as a plug-in estimator if an oracle reveals all information about $\bm{E}$ and $\bm{P}^{\star}$. 
In the sequel, our proof consists of two main steps: 
%
\begin{itemize}
	\item Show that the surrogate $\widetilde{v}_{i,j}$ is a reliable estimate of the truth ${v}^{\star}_{i,j}$, namely, $\widetilde{v}_{i,j}\approx {v}^{\star}_{i,j}$. 
	\item Show that the estimator in use and the surrogate estimator are sufficiently close, namely, $\widehat{v}_{i,j}\approx \widetilde{v}_{i,j}$.
\end{itemize}


\subsubsection{Step 1: show that $\widetilde{v}_{i,j}\approx {v}^{\star}_{i,j}$} 
%
Firstly, the fact that the $E_{i,j}$'s are zero-mean random variables
immediately reveals that $\widetilde{v}_{i,j}$ is an unbiased estimate
of $v_{i,j}^{\star}$, that is, 
%
\[
\mathbb{E}\big[\widetilde{v}_{i,j}\big]=v_{i,j}^{\star}.
\]
%
Secondly, given that the $E_{i,j}$'s are statistically independent,
we intend to invoke the Bernstein inequality to control the difference $\widetilde{v}_{i,j}-v_{i,j}^{\star}=\widetilde{v}_{i,j}-\mathbb{E}\big[\widetilde{v}_{i,j}\big]$.
To do so, one first calculates that
%
\begin{align*}
L_{0} & \coloneqq\max\left\{ \max_{l}E_{i,l}^{2}P_{l,j}^{\star 2},\,\max_{l}E_{l,j}^{2}P_{i,l}^{\star 2},\,E_{i,j}^{2}(P_{i,i}^\star +P_{j,j}^\star )^{2}\right\} \leq4B^{2}\|\bm{P}^\star \|_{\infty}^{2}
\end{align*}
%
and
\begin{align*}
V_{0} & \coloneqq\sum_{l:\,l\neq j}\mathsf{Var}\big(E_{i,l}^{2}\big)P_{l,j}^{\star 4}+\sum_{l:\,l\neq i}\mathsf{Var}\big(E_{l,j}^{2}\big)P_{i,l}^{\star 4}+\mathsf{Var}\big(E_{i,j}^{2}\big)(P_{i,i}^\star +P_{j,j}^\star )^{4}\\
 & \lesssim\sum_{l=1}^{n}\mathbb{E}\big[E_{i,l}^{4}\big]P_{l,j}^{\star 4}+\sum_{l=1}^{n}\mathbb{E}\big[E_{l,j}^{4}\big]P_{i,l}^{\star 4}+\mathbb{E}\big[E_{i,j}^{4}\big]P_{i,i}^{\star 2}P_{j,j}^{\star 2}\\
 & \lesssim B^{2}\|\bm{P}^\star \|_{\infty}^{2}\left\{ \sum_{l=1}^{n}\mathbb{E}\big[E_{i,l}^{2}\big]P_{l,j}^{2}+\sum_{l=1}^{n}\mathbb{E}\big[E_{l,j}^{2}\big]P_{i,l}^{\star 2}+\mathbb{E}\big[E_{i,j}^{2}\big]P_{i,i}^\star P_{j,j}^\star \right\} \\
 & \lesssim \sigma^{2} B^{2} \|\bm{P}^\star \|_{\infty}^{2}\left\{ \sum_{l=1}^{n}P_{l,j}^{\star 2}+\sum_{l=1}^{n}P_{i,l}^{\star 2}\right\} \\
 & \lesssim \sigma^{2} B^{2} \|\bm{P}^\star \|_{\infty}^{3},
\end{align*}
%
where the last line follows from \eqref{eq:P-jl-square-sum}. 
%since, for any $1\leq j\leq n$, 
%%
%\[
%\sum_{l=1}^{n}P_{l,j}^{2}=\big\|\bm{U}^{\star}\big(\bm{U}_{j,\cdot}^{\star}\big)^{\top}\big\|_{2}^{2}=\big\|\bm{U}_{j,\cdot}^{\star}\big\|_{2}^{2}=P_{j,j}\leq\|\bm{P}\|_{\infty}.
%\]
%
Invoking the Bernstein inequality (cf.~Corollary~\ref{thm:matrix-Bernstein-friendly}) reveals that with probability exceeding $1-O(n^{-5})$, 
%
\begin{align}
\big|\widetilde{v}_{i,j}-v_{i,j}^{\star}\big| & =\left|\widetilde{v}_{i,j}-\mathbb{E}\left[\widetilde{v}_{i,j}\right]\right|\lesssim\sqrt{V_{0}\log n}+L_{0}\log n \notag\\
 & \lesssim\sigma B\sqrt{\|\bm{P}^\star \|_{\infty}^{3}\log n}+B^{2}\|\bm{P}^\star \|_{\infty}^{2}\log n \notag\\
 & \lesssim\frac{\sigma B\mu^{3/2}r^{3/2}\sqrt{\log n}}{n^{3/2}}+\frac{\mu^{2}r^{2}B^{2}\log n}{n^{2}} ,
	\label{eq:tilde-v-star-v-diff}
\end{align}
%
where the last inequality results from \eqref{eq:P-inf-norm-UB}. 



\subsubsection{Step 2: show that $\widehat{v}_{i,j}\approx \widetilde{v}_{i,j}$} 
%
In order to accomplish this, we are in need of controlling the difference between $E_{i,j}$ (resp.~$P^\star_{i,j}$) and $\widehat{E}_{i,j}$ (resp.~$\widehat{P}_{i,j}$). 
To this end, apply the entrywise estimation guarantees in Corollary~\ref{cor:entrywise-error-general} to yield
%
\begin{align}
\big\|\widehat{\bm{E}}-\bm{E}\big\|_{\infty} & =\big\|\big(\bm{M}-\bm{U}\bm{\Lambda}\bm{U}^{\top}\big)-\big(\bm{M}-\bm{M}^{\star}\big)\big\|_{\infty} \notag\\
 & =\big\|\bm{U}\bm{\Lambda}\bm{U}^{\top}-\bm{M}^{\star}\big\|_{\infty}\lesssim\sigma\kappa^{2}\mu r\sqrt{\frac{\log n}{n}},
	\label{eq:entrywise-estimation-E}
\end{align}
%
and as a result,
%
\begin{equation}
\big\|\widehat{\bm{E}}\big\|_{\infty}\leq\big\|\bm{E}\big\|_{\infty}+\big\|\widehat{\bm{E}}-\bm{E}\big\|_{\infty}
	\lesssim B+\sigma\kappa^{2}\mu r\sqrt{\frac{\log n}{n}}
	\asymp B,
	\label{eq:entrywise-estimation-E-1}
\end{equation}
%
provided that $B\gtrsim \sigma\kappa^{2}\mu r\sqrt{\frac{\log n}{n}}$ (which is trivially satisfied if $\kappa^{2}\mu r\sqrt{\frac{\log n}{n}}\leq 1$). Moving on to the error term $\widehat{P}_{i,j}-P_{i,j}^\star $, we observe that
%
\begin{align*}
 & \big\|\widehat{\bm{P}}-\bm{P}^\star \big\|_{\infty} =\big\|\bm{U}\bm{U}^{\top}-\bm{U}^{\star}\bm{U}^{\star\top}\big\|_{\infty}\\
 & \qquad \leq\big\|\big(\bm{U}\mathsf{sgn}(\bm{H})-\bm{U}^{\star}\big)\bm{U}^{\star\top}\big\|_{\infty}+\big\|\bm{U}\mathsf{sgn}(\bm{H})\big(\bm{U}\mathsf{sgn}(\bm{H})-\bm{U}^{\star}\big)^{\top}\big\|_{\infty}\\
 & \qquad \leq\big\|\bm{U}\mathsf{sgn}(\bm{H})-\bm{U}^{\star}\big\|_{2,\infty}\left\{ \big\|\bm{U}^{\star}\big\|_{2,\infty}+\big\|\bm{U}\mathsf{sgn}(\bm{H})\big\|_{2,\infty}\right\} \\	
 & \qquad \leq\big\|\bm{U}\mathsf{sgn}(\bm{H})-\bm{U}^{\star}\big\|_{2,\infty}\left\{ \big\|\bm{U}^{\star}\big\|_{2,\infty}+\big\|\bm{U}\big\|_{2,\infty}\right\} .
\end{align*}
%
Taking this  together with the bound \eqref{eq:UsgnH-Ustar-bound-theorem-general} in Theorem~\ref{thm:UsgnH-Ustar-MUstar-general}, 
the incoherence assumption, and the inequality \eqref{eq:incoherence-U-estimate}, we arrive at
%
\begin{align}
	\big\|\widehat{\bm{P}}-\bm{P}^\star \big\|_{\infty} & \lesssim \frac{(\sigma\kappa\sqrt{\mu r}+\sigma\sqrt{r\log n})}{|\lambda_{r}^{\star}|} \sqrt{\frac{\mu r}{n}}. \label{eq:P-hat-Linf-norm-diff}
\end{align}
%
This taken together with \eqref{eq:P-inf-norm-UB} indicates that
%
\begin{align}
\big\|\widehat{\bm{P}}\big\|_{\infty} & \leq\big\|\bm{P}^\star \big\|_{\infty}+\big\|\widehat{\bm{P}}-\bm{P}^\star \big\|_{\infty}\nonumber \\
 & \lesssim \frac{\mu r}{n}+\frac{(\sigma\kappa\sqrt{\mu r}+\sigma\sqrt{r\log n})}{|\lambda_{r}^{\star}|}\sqrt{\frac{\mu r}{n}}\asymp\frac{\mu r}{n},\label{eq:P-hat-Linf-norm}
\end{align}
%
with the proviso that $\sigma\sqrt{n}\lesssim|\lambda_{r}^{\star}|/\kappa$ and $\sigma\sqrt{n\log n}\lesssim|\lambda_{r}^{\star}|$. 



Armed with the preceding bounds, we are now positioned to control $\widehat{v}_{i,j} - \widetilde{v}_{i,j}$. From the definition of $\widetilde{v}_{i,j}$ and $v_{i,j}^{\star}$, we recognize that
%
\begin{align}
\big| \widehat{v}_{i,j} - \widetilde{v}_{i,j} \big| & \leq\underset{\eqqcolon\,\alpha_{1}}{\underbrace{\left|\sum_{l=1}^{n}\left(\widehat{E}_{i,l}^{2}\widehat{P}_{l,j}^{2}-E_{i,l}^{2}P_{l,j}^{\star 2}\right)\right|}}+\underset{\eqqcolon\,\alpha_{2}}{\underbrace{\left|\sum_{l=1}^{n}\left(\widehat{P}_{i,l}^{2}\widehat{E}_{l,j}^{2}-P_{i,l}^{\star 2}E_{l,j}^{2}\right)\right|}}\nonumber \\
 & \qquad+2\, \underset{\eqqcolon\,\alpha_{3}}{\underbrace{\left|\widehat{E}_{i,j}^{2}\widehat{P}_{i,i}\widehat{P}_{j,j}-E_{i,j}^{2}P_{i,i}^\star P_{j,j}^\star \right|}},
	\label{eq:v-tilde-bound-alpha-123}
\end{align}
%
leaving us with three terms to cope with. 
Regarding the first term $\alpha_1$ on the right-hand side of \eqref{eq:v-tilde-bound-alpha-123}, it can be easily verified that
%
\begin{align}
\alpha_{1} & \leq
\left|\sum_{l=1}^{n}\left(\widehat{E}_{i,l}^{2}\widehat{P}_{l,j}^{2}-E_{i,l}^{2}\widehat{P}_{l,j}^{2}\right)\right|+\left|\sum_{l=1}^{n}\left(E_{i,l}^{2}\widehat{P}_{l,j}^{2}-E_{i,l}^{2}P_{l,j}^{\star2}\right)\right| \notag\\
 & \leq \left(\max_{l}\left|\widehat{E}_{i,l}^{2}-E_{i,l}^{2}\right|\right)\sum_{l=1}^{n}\widehat{P}_{l,j}^{2}+\left(\max_{l}\left| \widehat{P}_{l,j}^{2} - P_{l,j}^{\star 2} \right|\right)\sum_{l=1}^{n}E_{i,l}^{2} \notag\\
	& \overset{\mathrm{(i)}}{\leq} \left(\|\bm{E}\|_{\infty}+\|\widehat{\bm{E}}\|_{\infty}\right)\big\|\widehat{\bm{E}}-\bm{E}\big\|_{\infty}\widehat{P}_{j,j} \notag\\
	&\qquad +\left(\|\bm{P}^\star \|_{\infty}+\|\widehat{\bm{P}}\|_{\infty}\right)\big\|\widehat{\bm{P}}-\bm{P}^\star \big\|_{\infty}\|\bm{E}\|^{2} \notag\\
	& \overset{\mathrm{(ii)}}{\lesssim} B\sigma\kappa^{2}\mu r\sqrt{\frac{\log n}{n}}\cdot\frac{\mu r}{n}+\frac{\mu r}{n}\cdot\frac{(\sigma\kappa\sqrt{\mu r}+\sigma\sqrt{r\log n})}{|\lambda_{r}^{\star}|}\sqrt{\frac{\mu r}{n}}\cdot\sigma^{2}n.  \notag
	%\label{eq:alpha1-ii-UB}
\end{align}
%
Here, (i) makes use of \eqref{eq:P-jl-square-sum} (with $\bm{P}^{\star}$ replaced by $\widehat{\bm{P}}$), 
whereas (ii) holds true due to \eqref{eq:entrywise-estimation-E}, \eqref{eq:entrywise-estimation-E-1}, \eqref{eq:P-inf-norm-UB}, \eqref{eq:P-hat-Linf-norm-diff}, \eqref{eq:P-hat-Linf-norm}, and Lemma~\ref{lemma:general-noise-bound}.  
The second term $\alpha_2$ on the right-hand side of \eqref{eq:v-tilde-bound-alpha-123} can be bounded in the same manner and we omit it here for brevity. 
When it comes to the last term $\alpha_3$ on the right-hand side of \eqref{eq:v-tilde-bound-alpha-123}, one has
%
\begin{align*}
\alpha_{3} & \leq\left|\widehat{E}_{i,j}^{2}-E_{i,j}^{2}\right|\left|\widehat{P}_{i,i}\widehat{P}_{j,j}\right|+ E_{i,j}^{2} \left|\widehat{P}_{i,i}\widehat{P}_{j,j}-P_{i,i}^\star P_{j,j}^\star \right|\\
 & \lesssim\big\|\widehat{\bm{E}}-\bm{E}\big\|_{\infty}\left\{ \big\|\bm{E}\big\|_{\infty}+\big\|\widehat{\bm{E}}\big\|_{\infty}\right\} \big\|\widehat{\bm{P}}\big\|_{\infty}^{2} \notag\\
	& \qquad +B^{2}\big\|\widehat{\bm{P}}- \bm{P}^{\star} \big\|_{\infty}\left\{ \big\|\bm{P}^\star \big\|_{\infty}+\big\|\widehat{\bm{P}}\big\|_{\infty}\right\} \\
 & \lesssim\sigma\kappa^{2}\mu r\sqrt{\frac{\log n}{n}}\cdot B\cdot\left(\frac{\mu r}{n}\right)^{2}+B^{2}\cdot\frac{(\sigma\kappa\sqrt{\mu r}+\sigma\sqrt{r\log n})}{|\lambda_{r}^{\star}|}\sqrt{\frac{\mu r}{n}}\cdot\frac{\mu r}{n} \\
 & \lesssim B\sigma\kappa^{2}\mu r\sqrt{\frac{\log n}{n}}\cdot\frac{\mu r}{n}+\frac{\mu r}{n}\cdot\frac{(\sigma\kappa\sqrt{\mu r}+\sigma\sqrt{r\log n})}{|\lambda_{r}^{\star}|}\sqrt{\frac{\mu r}{n}}\cdot\sigma^{2}n, 
\end{align*}
%
where the penultimate inequality is a consequence of \eqref{eq:entrywise-estimation-E}, \eqref{eq:entrywise-estimation-E-1}, \eqref{eq:P-inf-norm-UB}, \eqref{eq:P-hat-Linf-norm-diff} and \eqref{eq:P-hat-Linf-norm},
and the last line holds true as long as $\mu r \leq n$ (cf.~\eqref{eq:mu-constraint-general}) and $B\lesssim \sigma \sqrt{n}$. 
Combining the above inequalities allows one to reach
%
\begin{align*}
 & \big| \widehat{v}_{i,j} - \widetilde{v}_{i,j} \big|  \leq\alpha_{1}+\alpha_{2}+\alpha_{3}\\
 & \qquad \lesssim  
	 B\sigma\kappa^{2}\mu r\sqrt{\frac{\log n}{n}}\cdot\frac{\mu r}{n}+\frac{\mu r}{n}\cdot\frac{\sigma\kappa\sqrt{\mu r}+\sigma\sqrt{r\log n}}{|\lambda_{r}^{\star}|}\sqrt{\frac{\mu r}{n}}\cdot\sigma^{2}n \\
 & \qquad \asymp \frac{B\sigma\kappa^{2}\mu^{2}r^{2}\sqrt{\log n}}{n^{3/2}}+\frac{\sigma^{3}\kappa\mu^{2}r^{2}+\sigma^{3}\mu^{3/2}r^{2}\sqrt{\log n}}{|\lambda_{r}^{\star}|\sqrt{n}}
\end{align*}
%
with probability at least $1-O(n^{-5})$. 
%provided that $\mu r\leq n$. 




\subsubsection{Step 3: combining the above bounds} 

Putting together the results in the previous steps, we can readily derive
%
\begin{align}
\big|\widehat{v}_{i,j}-v_{i,j}^{\star}\big| & \leq\big|\widetilde{v}_{i,j}-v_{i,j}^{\star}\big|+ \big| \widehat{v}_{i,j} - \widetilde{v}_{i,j} \big|  \notag\\
 & \lesssim\left(\frac{\sigma B\mu^{3/2}r^{3/2}\sqrt{\log n}}{n^{3/2}}+\frac{\mu^{2}r^{2}B^{2}\log n}{n^{2}}\right) \notag\\
 & \qquad+\left(\frac{B\sigma\kappa^{2}\mu^{2}r^{2}\sqrt{\log n}}{n^{3/2}}+\frac{\sigma^{3}\kappa\mu^{2}r^{2}+\sigma^{3}\mu^{3/2}r^{2}\sqrt{\log n}}{|\lambda_{r}^{\star}|\sqrt{n}}\right) \notag\\
	& \asymp \frac{B\sigma\kappa^{2}\mu^{2}r^{2}\sqrt{\log n}}{n^{3/2}}+\frac{\sigma^{3}\kappa\mu^{2}r^{2}\sqrt{\log n}}{|\lambda_{r}^{\star}|\sqrt{n}},  \label{eq:vij-vhat-ij-gap-bound}
 %& \lesssim\frac{1}{\sqrt{\log n}}v_{i,j}^{\star},
\end{align}
%
where the last relation is guaranteed as long as $B\lesssim \sigma \sqrt{n/\log n}$. 
Consequently, if Condition~\eqref{eq:Uj-Ui-lower-bound-lem} holds, 
%$\frac{B\sigma\kappa^{2}\mu^{2}r^{2}\log n}{n^{3/2}}+\frac{\sigma^{3}\kappa\mu^{2}r^{2}\log n}{|\lambda_{r}^{\star}|\sqrt{n}}
%\lesssim\sigma_{\min} ^2 \big\{ \|\bm{U}_{j,\cdot}^{\star}\|_{2}^{2}+\|\bm{U}_{i,\cdot}^{\star}\|_{2}^{2} \big\} $, 
then it follows from \eqref{eq:vij-vhat-ij-gap-bound} and the lower bound \eqref{eq:vij-lower-bound-12} that
%
\begin{align*}
	\big|\widehat{v}_{i,j}-v_{i,j}^{\star}\big| \lesssim \frac{1}{\sqrt{\log n}} v_{i,j}^{\star}, 
\end{align*}
%
thus concluding the proof. 





